% !TeX root = main.tex 
%!TEX root = main.tex

\newcommand{\maxwalk}[1]{\ensuremath{\textnormal{\textsf{w}}_{\textnormal{max}}(#1)}}
\newcommand{\timemax}[1]{\ensuremath{\textnormal{\textsf{t}}_{\textnormal{max}}(#1)}}

\newcommand{\maxwalkW}{\ensuremath{\textnormal{\textsf{w}}_{\textnormal{max}}}}
\newcommand{\timemaxW}{\ensuremath{\textnormal{\textsf{t}}_{\textnormal{max}}}}

\newcommand{\dist}{\ensuremath{\mathcal{D}}}



\section{Preliminaries}\label{sec:prelim}

\paragraph{Notation.} For integers $1 \leq a \leq b$, we define $[a:b] := \set{a,a+1,\ldots,b}$ and $[b] := \set{1,\ldots,b}$. For a tuple $x=(x_1,\ldots,x_t)$ and $i \in [2:t]$, 
we denote $x_{<i} := (x_1,\ldots,x_{i-1})$.  

For a random variable $X$, we sometimes write the subscript $X$ in $\Pr_X(f(X))$ (for some $f$), to 
explicitly specify the randomness is coming from $X$.  
For a random variable $X$ and a set $S$, we write $X \in_R S$ to mean $X$ is sampled uniformly at random from $S$. 

\paragraph{Gambler's ruin.} We use a standard type of random walks on integer line $\set{0,1,\ldots,2m}$ for some integer $m \geq 1$, often called the \emph{gambler's ruin}; see, e.g.~\cite[Section 2.1]{levin2026markov}. 

The walk $W=w_1,w_2,\cdots$ starts at some integer $w_1 \in [2m-1]$ and in each time $t \geq 1$, it updates $w_{t+1} = w_{t}+1$ with probability half and $w_{t+1} = w_{t}-1$ otherwise. 
The walk terminates at the first time $t$ when $w_t = 0$ or $w_t = 2m$ and remains in this state. Throughout, by a random walk, we always mean the above type of walk unless explicitly stated otherwise. 

\begin{fact}[{c.f.~\cite[Proposition 2.1]{levin2026markov}}]\label{fact:ruin}
	For a random walk $W$ over $\set{0, \ldots, 2m}$ and an integer $k$ in $\set{0,\ldots,2m}$,
	\[
		\Pr_{W}\paren{\text{$W$ reaches $0$ before $2m$} \mid \text{$W$ starts at $k$}} = \frac{2m-k}{2m}. 
	\]
	Moreover, the expected length of the walk before terminating is at most $m^2$. 
\end{fact}

For a (finite) random walk $W = w_1,\cdots,w_T$, where $w_T \in \set{0,2m}$ is a terminating state, define: 
\begin{align*}
	\maxwalkW &= \maxwalk{W} := \max_{t \in [T]} \set{w_t} \quad \text{and}\quad 
	\timemaxW = \timemax{W} := \min \set{t \in [T] : w_t=\maxwalk{W}},  
\end{align*}
as the \emph{largest} state seen on the walk and the \emph{first} time the walk sees its eventual largest state. 

A corollary of~\Cref{fact:ruin} is the following. 
\begin{claim}\label{clm:ruin}
	For a random walk $W$ over $\set{0,\ldots,2m}$ and any $1 \leq i \leq j < 2m$,
	\[
		\Pr_W\paren{\text{$W$ ends in $0$} \wedge \maxwalk{W} = j \mid \text{$W$ starts at $i$}} = \frac{i}{j \cdot (j+1)}. 
	\]
\end{claim}
\begin{proof}
	For this to happen, we need $W$ to reach $j$ before $0$, and once on $j$, to reach $0$ before $j+1$. Given the random walk is memoryless, these two events are independent and thus we have
	\begin{align*}
		&\Pr_W\paren{\text{$W$ ends in $0$} \wedge \maxwalk{W} = j \mid \text{$W$ starts at $i$}} = \\
		&\hspace{15pt} \Pr_W\paren{\text{$W$ reaches $j$ before $0$} \mid \text{$W$ starts at $i$}} \cdot  \Pr_W\paren{\text{$W$ reaches $0$ before $j+1$} \mid \text{$W$ starts at $j$}}. 
	\end{align*}
	The first term is about a random walk on $\set{0,\ldots,j}$ starting from $i$ and thus by~\Cref{fact:ruin} has probability $i/j$. The second term is about a random walk on $\set{0,\ldots,j+1}$ starting from $j$
	and thus by~\Cref{fact:ruin} has probability $1/(j+1)$. Multiplying these two concludes the proof. 
\end{proof}
We also have the following standard claim on the length of random walks. 
\begin{claim}\label{clm:walk-length}
	The probability that a random walk $W$ starting from any state does not terminate within $2m^3$ steps is at most $2^{-m}$. 
\end{claim}
\begin{proof}
	The expected time for a walk $W$ to terminate, starting from any state, is at most $m^2$ by~\Cref{fact:ruin}. 
	Thus, by Markov bound, the probability that the walk has not terminated after $2m^2$ steps is at most half. Repeating the argument from current state, 
	the probability that $m$ repetitions, each of length $2m^2$, does not terminate the walk is at most $2^{-m}$. 
\end{proof}
